\chapter{Animal Adaptations}\label{ch:g4:adaptations}

\omterm{def:g1:plants-animals:animal}{Animals} live in \omterm{def:g2:habitats:kinds}{deserts}, \omterm{def:g2:habitats:kinds}{oceans}, \omterm{def:g2:habitats:kinds}{forests} and cities. Features that help
them survive and \omterm{def:g2:what-alive:signs}{reproduce} in their surroundings are called
\omterm{def:g4:adaptations:def}{adaptations}. This chapter explores \omterm{def:g1:my-body:body}{body} structures and behaviours that
match the challenges of different \omterm{def:g1:caring:needs}{places}.

\section{What is an adaptation?}

\begin{definition}[Adaptation]\label{def:g4:adaptations:def}
An \emph{adaptation}\index{adaptation} is a feature of a \omterm{def:g1:living-nonliving:living}{living thing}'s
\omterm{def:g1:my-body:body}{body} or behaviour that helps it survive and \omterm{def:g2:what-alive:signs}{reproduce} in its
environment. Adaptations are inherited over generations; they are not
choices made by one \omterm{def:g1:plants-animals:animal}{animal} in a single day.
\end{definition}

\begin{example}[Quick pictures]\label{ex:g4:adaptations:quick}
A duck's webbed \omterm{def:g1:my-body:parts}{feet} help it swim. A camel's ability to go a long time
between drinks suits dry regions. A hedgehog's spines discourage some
predators.
\end{example}

\section{Structural adaptations}

\begin{definition}[Structural adaptation]\label{def:g4:adaptations:struct}
A \emph{structural adaptation}\index{adaptation!structural} is a \omterm{def:g1:my-body:body}{body}
part or \omterm{def:g1:my-body:body}{body} feature that helps survival --- for example thick fur,
sharp claws, camouflage colour, or a streamlined shape.
\end{definition}

\begin{omfigure}
\begin{tikzpicture}[font=\small]
  % cold climate
  \draw[thick, omDef, fill=omDef!10] (0,0) ellipse (1.3 and 0.9);
  \fill[omDef!40] (-0.35,0.2) circle [radius=0.12];
  \fill[omDef!40] (0.35,0.2) circle [radius=0.12];
  \draw[thick, omDef] (-0.9,-0.5) -- (-1.1,-1.3);
  \draw[thick, omDef] (-0.4,-0.7) -- (-0.5,-1.35);
  \draw[thick, omDef] (0.4,-0.7) -- (0.5,-1.35);
  \draw[thick, omDef] (0.9,-0.5) -- (1.1,-1.3);
  \draw[thick, omDef, fill=omDef!20] (-1.5,0.2) ellipse (0.35 and 0.2);
  \node at (0,-1.8) {thick fur, small ears};
  \node[omDef] at (0,1.3) {cold place};
  % hot climate
  \begin{scope}[xshift=5.2cm]
    \draw[thick, omProp, fill=omProp!12] (0,0) ellipse (1.1 and 0.75);
    \fill[omProp!50] (-0.3,0.15) circle [radius=0.1];
    \fill[omProp!50] (0.3,0.15) circle [radius=0.1];
    \draw[thick, omProp, fill=omProp!15] (-0.5,0.9) ellipse (0.15 and 0.45);
    \draw[thick, omProp, fill=omProp!15] (0.5,0.9) ellipse (0.15 and 0.45);
    \draw[thick, omProp] (-0.7,-0.5) -- (-0.8,-1.2);
    \draw[thick, omProp] (0.7,-0.5) -- (0.8,-1.2);
    \node at (0,-1.8) {large ears, thin coat};
    \node[omProp] at (0,1.55) {hot place};
  \end{scope}
\end{tikzpicture}

{\small \omterm{def:g1:my-body:body}{Body} features often match climate: insulation and reduced heat
loss in the cold; heat release and \omterm{def:g4:plants-need:needs}{light} coats in the heat (schematic).}
\end{omfigure}

\begin{example}[Camouflage]\label{ex:g4:adaptations:camo}
A stick \omterm{def:g3:insects:insect}{insect}'s shape and colour hide it among twigs. A flounder's
flat \omterm{def:g1:my-body:body}{body} and mottled \omterm{def:g1:senses:five}{skin} blend with the sea floor. Camouflage can
protect prey --- or help predators sneak up.
\end{example}

\section{Behavioural adaptations}

\begin{definition}[Behavioural adaptation]\label{def:g4:adaptations:behav}
A \emph{behavioural adaptation}\index{adaptation!behavioural} is a
helpful way of acting: migrating, hibernating, hunting at night,
building a nest, or living in a group.
\end{definition}

\begin{example}[Seasonal moves]\label{ex:g4:adaptations:migrate}
Many \omterm{def:g2:animal-groups:five}{birds} migrate to find \omterm{def:g2:food-energy:food}{food} and nesting sites as \omterm{def:g1:seasons:season}{seasons} change.
Some \omterm{def:g2:animal-groups:five}{mammals} sleep through the hardest \omterm{def:g1:seasons:four}{winter} months. These behaviours
match changing conditions.
\end{example}

\begin{proposition}[Structure and behaviour together]\label{prop:g4:adaptations:both}
Survival often depends on both \omterm{def:g1:my-body:body}{body} features and behaviour. A \omterm{def:g2:animal-groups:five}{bird} may
have wings (structure) and a migration route (behaviour).
\end{proposition}
\admitted

\section{Matching habitat}

\omimg{camouflage.jpg}{0.52\linewidth}{Different \omterm{def:g2:habitats:habitat}{habitats} pose different
problems: finding \omterm{def:g4:plants-need:needs}{water}, staying warm, escaping predators, catching
\omterm{def:g2:food-energy:food}{food}.}

\begin{method}[Reading an animal's lifestyle]\label{met:g4:adaptations:read}
\begin{enumerate}
  \item Name the \omterm{def:g2:habitats:habitat}{habitat} (\omterm{def:g2:habitats:kinds}{desert}, \omterm{def:g2:habitats:kinds}{ocean}, \omterm{def:g2:habitats:kinds}{forest}, \omterm{def:g2:habitats:kinds}{grassland}, \ldots).
  \item List challenges (heat, cold, scarcity of \omterm{def:g4:plants-need:needs}{water}, hiding, speed).
  \item Link each major \omterm{def:g1:my-body:body}{body} feature or behaviour to a challenge.
  \item Ask: how would this \omterm{def:g1:plants-animals:animal}{animal} struggle in a very different \omterm{def:g1:caring:needs}{place}?
\end{enumerate}
\end{method}

\begin{example}[Polar bear sketch]\label{ex:g4:adaptations:polar}
Thick fur and fat keep heat in; large paws spread weight on snow and
ice; white coat provides camouflage; hunting seals on ice is a
behaviour matched to the Arctic \omterm{def:g2:food-energy:food}{food} supply.
\end{example}

\begin{omfigure}
\begin{tikzpicture}[font=\footnotesize,
  box/.style={draw, rounded corners, minimum width=2.4cm, minimum height=0.75cm,
    align=center, fill=omDef!10}]
  \node[box] (h) at (0,1.6) {habitat challenge};
  \node[box, fill=omMeth!12] (s) at (-2.6,0) {structural\\adaptation};
  \node[box, fill=omProp!12] (b) at (2.6,0) {behavioural\\adaptation};
  \node[box, fill=omThm!12] (r) at (0,-1.6) {better chance to\\survive and reproduce};
  \draw[-Stealth, thick] (h) -- (s);
  \draw[-Stealth, thick] (h) -- (b);
  \draw[-Stealth, thick] (s) -- (r);
  \draw[-Stealth, thick] (b) -- (r);
\end{tikzpicture}

{\small Challenges of a \omterm{def:g1:caring:needs}{place} shape which features help --- in \omterm{def:g1:my-body:body}{body} and
in behaviour.}
\end{omfigure}

\section{Humans and tools}

\section{Exercises}

\begin{exercise}[$\star$]\label{exo:g4:adaptations:1}
What is an \omterm{def:g4:adaptations:def}{adaptation}?
\end{exercise}

\begin{exercise}[$\star$]\label{exo:g4:adaptations:2}
Give one \omterm{def:g4:adaptations:struct}{structural adaptation} and one \omterm{def:g4:adaptations:behav}{behavioural adaptation} of any
\omterm{def:g1:plants-animals:animal}{animal} you know.
\end{exercise}

\begin{exercise}[$\star$]\label{exo:g4:adaptations:3}
How can camouflage help either a prey \omterm{def:g1:plants-animals:animal}{animal} or a predator?
\end{exercise}

\begin{exercise}[$\star$]\label{exo:g4:adaptations:4}
Why is thick fur helpful in a cold \omterm{def:g2:habitats:habitat}{habitat}?
\end{exercise}

\begin{exercise}[$\star\star$]\label{exo:g4:adaptations:5}
A \omterm{def:g2:habitats:kinds}{desert} fox has large \omterm{def:g1:senses:five}{ears}. Suggest how this feature might help in a
hot, dry \omterm{def:g1:caring:needs}{place}.
\end{exercise}

